\section{系統模型}
\label{sec:model}

\subsection{架構}

本文建模一個單層分散式服務：$N$ 個同質工作節點組成叢集，前端設有一個負載均衡器。
請求依泊松過程抵達系統，
總到達速率為 $\Lambda = L \cdot N \cdot \mu$，
其中 $L \in (0, 1)$ 為\emph{負載因子}，
$\mu = 1/\tbase$ 為單一節點之服務速率。
各工作節點以先進先出方式逐一處理請求（單伺服器模型）。
正常情況下服務時間為定值 $\tbase$，整體符合 M/D/1 佇列模型。

負載均衡器採用加權 Power-of-Two-Choices（P2C）策略：
針對每筆請求，依權重 $w_i \in [0, 1]$ 按比例隨機抽取兩個候選節點，
再將請求路由至其中佇列較短者。

\subsection{故障模型}

當節點 $i$ 發生\emph{慢故障}時，
其服務時間將乘以減速因子 $s_i > 1$：
\[
  t_i = s_i \cdot \tbase
\]
假設 $N$ 個節點中同時有 $f$ 個發生故障，
且故障具有\emph{持續性}：
一旦發生，減速效果將延續至系統主動恢復為止（若有恢復機制）。
本文探討以下五種故障模式：

\begin{itemize}
\item \textbf{永久性}：故障發生後減速因子 $s$ 維持恆定。
\item \textbf{漸進式}：$s(t) = 1 + (s_{\max}-1)(1-e^{-\beta t})$，
  減速隨時間逐漸加劇。
\item \textbf{波動式}：在 $s = s_{\text{peak}}$ 和 $s = 1$ 之間交替，
  具有可配置的開／關持續時間。
\item \textbf{級聯式}：
  故障依交錯的時間間隔先後擴散至不同節點。
\item \textbf{恢復式}：故障活躍一段固定時間後，節點恢復正常。
\end{itemize}

\subsection{效能指標}

本文以下列尾端延遲指標評估系統效能：
\begin{itemize}
\item \textbf{P99/P999 延遲}：
  請求端到端延遲之第 99 及第 99.9 百分位（含排隊等候時間）。
\item \textbf{SLO 違反率}：
  延遲超過 SLO 閾值（實驗中設為 50\,ms）之請求佔比。
\item \textbf{有效吞吐量}（goodput）：每秒在 SLO 內完成的請求數。
\item \textbf{受影響比例}：
  經由故障節點處理之請求佔比，用以衡量負載均衡器的分流效果。
\end{itemize}

\subsection{極小極大遺憾}

為比較各策略在多種場景下的效能，
本文採用\emph{極小極大遺憾}（minimax regret）準則。
給定策略 $j$ 與場景 $i$，其遺憾定義為：
\[
  R_{ij} = P99_{ij} - \min_k P99_{ik}
\]
策略 $j$ 的極小極大遺憾為 $\max_i R_{ij}$，
即該策略在最不利場景下相對於最優策略的效能落差。
極小極大遺憾愈低，代表策略在各種運行條件下愈具穩健性。
